\documentclass[11pt, a4paper]{article}
\input{preamble.tex}
\switchlang{ru}
\begin{document}

\title{1981 "--- Xerox Star Mechanical Mouse}
\date{}
\maketitle
\selectlanguage{russian}

В апреле 1981 года корпорация Xerox анонсировала Xerox Star 8010 Information System, представлявшую собой первый офисный компьютер, а также первую многооконную графическую оболочку на основе иконок и папок и первую двухкнопочную мышь \cite{star}. Мышь выпускалась в двух вариантах: оптическая мышь Xerox Star, изобретенная специально для этих компьютеров, и механическая мышь. Механическая мышь представляла собой разновидность более старого варианта мыши компьютеров Xerox Alto с конструкцией на основе шара, разработанной Джеком Хоули (одним из авторов изначального патента Xerox 1973 года на мышь с двумя наклонными колесами \cite{pat}). Оба варианта мышей Xerox Star выпускались с одинаковой верхней частью корпуса (рис. \ref{fig:XeroxStarMechanicalPic}), представлявшей собой дальнейшее эволюционное развитие визуального стиля мыши Xerox Alto (первой массово производившейся мыши в истории компьютеров). Вероятнее всего, механическая мышь использовалась в основном в начальном периоде существования Xerox Star, когда оптические мыши еще не производились в достаточном количестве: сама по себе конструкция Хоули, согласно цитатам из \cite{mouses}, была не слишком надежной: она собирала грязь, быстро засорядалсь и переставала управлять курсором, что, собственно, и стало основной причиной изобретения оптической альтернативы. Тем не менее, в документации Xerox механические мыши продолжали упоминаться как альтернатива до самого конца выпуска рабочих станций Star 8010 в 1985 году, а пришедшие им на смену компьютеры Xerox 6085 хотя и оснащались исключительно оптической мышью с разъемом нового типа, но по какой-то не вполне очевидной причине дополнительно комплектовались еще и ковриком для механической мыши \cite[стр. 133]{aimouse}.

\begin{figure}[h]
    \centering
    \includegraphics[scale=0.7]{1981_xerox_star_mechanical_mouse/pic_30.jpg}
    \caption{Xerox Star Mechanical Mouse}
    \label{fig:XeroxStarMechanicalPic}
\end{figure}

\section*{Внешний вид}

Мышь сохраняет преемственность с оригинальной прямоугольной формой мыши Alto, но в более элегантном исполнении. Верхняя часть, выполненная из бежевого пластика, по-прежнему напоминает параллелепипед. Однако она расширяется к низу, её передняя и задняя стороны выпуклые и заметно наклонены, а вертикальные ребра закруглены, что существенно уменьшает первоначальную ассоциацию с «коробкой». На верхней стороне корпуса расположены две большие продолговатые кнопки из более темного пластика, которые в основном повторяют форму и размер кнопок мыши Alto, но меньше закруглены (рис. \ref{XeroxStarMechanicalTopAndBottom}). Поскольку добавлять третью кнопку между ними не планировалось, две кнопки расположены ближе друг к другу. Стоит отметить, что при создании более поздней модели, Xerox 6085, дизайнерам мыши пришлось отодвинуть кнопки друг от друга, чтобы зарезервировать место для трехкнопочного варианта.

\begin{figure}[h]
    \centering
    \includegraphics[scale=0.52]{1981_xerox_star_mechanical_mouse/top_30.jpg}
    \includegraphics[scale=0.52]{1981_xerox_star_mechanical_mouse/bottom_30.jpg}
    \caption{Xerox Star Mechanical Mouse, вид сверху и снизу}
    \label{fig:XeroxStarMechanicalTopAndBottom}
\end{figure}

Нижняя стоорона целиком выполнена из металла (рис. \ref{fig:XeroxStarMechanicalTopAndBottom}). Вращение регистрируется гладким стальным шаром в центре, а еще два шарика меньшего размера играют роль ножек для минимизации трения. Съемное кольцо, позволяющее извлечь шар для удаления собравшегося мусора, в данной модели еще не предусмотрено, поэтому для чистки необходима полная разборка.

\section*{Размеры}

Корпус механической мыши Xerox Star содержит те же внутренние компоненты, что и мышь Alto, поэтому новая форма с наклонными боковыми сторонами вынудила дизайнеров немного увеличить общие габариты в нижней части корпуса. Теоретически, оптическая мышь Xerox Star могла бы быть меньше и ниже, но дизайнеры не пошли по этому пути ради достижения максимально близкого пользовательского опыта для всех конфигураций Xerox Star. Тем не менее, по меркам 1980-х годов мышь Xerox Star можно назвать скорее компактной, чем средних размеров, а визуально она выглядит еще меньшей (рис. \ref{fig:XeroxStarMechanicalSize}).

\begin{figure}[h]
    \centering
    \includegraphics[scale=0.45]{1981_xerox_star_mechanical_mouse/size_30.jpg}
    \caption{Xerox Star Mechanical Mouse на размерном коврике с шагом сетки 1~см}
    \label{fig:XeroxStarMechanicalSize}
\end{figure}

\section*{Эргономика}

В плане эргономики и во внешнем виде Xerox Star Mouse прослеживается минимализм. Удобство использования, безусловно, меньше страдает от «прямоугольной» формы корпуса, чем это было с мышью Alto. Выпуклые продолговатые кнопки, расположенные в пределах досягаемости пальцев, по-прежнему удобны, а наклон ближней к пользователю стороны корпуса даже обеспечивает некоторый комфорт для ладони (рис. \ref{fig:XeroxStarMechanicalHand}) "--- по крайней мере, больше, чем у многих мышей 1980-х годов (конечно, более низкий корпус сделал бы мышь еще удобнее). Кроме того, узкий высокий корпус позволяет переносить мышь на другую часть коврика, удерживая ее с боков пальцами. Это может оказаться полезным для преодоления временных застреваний блока механического энкодера из-за попавшего внутрь фрагмента мусора "--- но опять же, не настолько, учитывая, что устройство достаточно тяжелое из-за металлических деталей.

\begin{figure}[h]
    \centering
    \includegraphics[scale=0.50]{1981_xerox_star_mechanical_mouse/hand_30.jpg}
    \caption{Xerox Star Mechanical Mouse с моделью руки человека}
    \label{fig:XeroxStarMechanicalHand}
\end{figure}

\section*{Подключение}

Как и все мыши Xerox, механическая мышь Xerox Star имеет квадратурный интерфейс. Разъем подключения такой же как у мыши Xerox Alto (в данном варианте "--- DA-15) и имеет то же расположение контактов \cite{pinout}.

\section*{Внутреннее устройство}

В разобранном виде манипулятор показан на рис. \ref{fig:XeroxStarMechanicalInside}. Можно отметить съемную глухую защиту шара, требующую дополнительных операций разборки для удаления мусора. В мыши испольованы механические энкодеры, представляющие собой металлический контактный барабан (вместо более распространенного в более поздних моделях диска) и четыре контакта различной длины вместо двух для увеличения разрешающей способности энкодера.

\begin{figure}[h]
    \centering
    \includegraphics[width=\textwidth]{1981_xerox_star_mechanical_mouse/inside_30.jpg}
    \caption{Xerox Star Mechanical Mouse в разобранном виде}
    \label{fig:XeroxStarMechanicalInside}
\end{figure}

\begin{thebibliography}{9}
\bibitem {star} Xerox Star: Wikipedia \url{https://en.wikipedia.org/wiki/Xerox_Star}
\bibitem{pat} Transducer for a display-oriented pointing device \url{https://patents.google.com/patent/US3892963A/en}
\bibitem{mouses} Xerox Mice. oldmouse.com \url{https://web.archive.org/web/20210418000634/http://oldmouse.com/mouse/xerox/alto.shtml}
\bibitem{aimouse} Xerox 6085 PCS Service Manual. October 1988 \url{https://www.bitsavers.org/pdf/xerox/6085/service/700P05710_6085_PCS_Service_Manual_Oct88.pdf} P. 125
\bibitem{pinout} 10-Pin IDC header in Lyon mouse -- bitsavers.org \url{https://www.bitsavers.org/pdf/xerox/mouse/lyon_optical/Lyon_Mouse_Notes.txt}
\end{thebibliography}
\end{document}
